\scanpage{45}

dictionnaire, $\Gamma$ s'interprète comme le \deleted{\unsure{tétr}} Groupe des autom du tétraèdre, ou du cube orienté, $\Gamma^{+}$ comme le gpe des autom \uline{orientés} du tétraèdre (ou celui des automorphismes ``biorientés'' du cube) \ldots

2. On s'intéresse \deleted{plus \unsure{m}} aux systèmes $\deleted{\unsure{(H\subset G,)}}\inserted{(\mathbb A,\deleted{\unsure{G}}\mathbb H,}\rho,\rho',\rho'')$ où $\rho,\rho',\rho''$ sont des sections \deleted{qui} \inserted{de} $G$, \inserted{telles que $\rho\rho'\rho''=1$} qui engendrent un ss-gpe fini $\mathbb H$, non abélien ni diédral (i.e. $\simeq \unsure{A_4|\mathcal A_4|\mathfrak A_4},\unsure{\mathfrak S_4|\mathcal S_4}$ ou $A_5$) -- Donc $H$ sera une forme splittée de l'un de ces gpes (mais $X$ \inserted{ou $G$} ne sera pas splitté pour autant --)

% slanted
\margin{\uline{NB} $\mathbb H$ résulte de $\rho,\rho',\rho''$, et $\rho''$ résulte de $\rho,\rho'$ (et $\rho$ de $\rho',\rho''$, $\rho'$ de $\rho'',\rho\ldots$)}

\unsure{On} \unsure{peut} dire aussi (sur base $S$ connexe) qu'à la donnée $\rho,\rho',\rho''$ correspond un tétraèdre \inserted{comb.} \unsure{resp.} un cube orienté \inserted{comb.} (\unsure{resp|ou} un biicosaèdre comb., cf plus bas) qui définit $X$ (modulo grain de sel du cas icosaédral \ldots), qu'on peut oublier. Et on se propose de trouver les syst. de trois gén. $\rho,\rho',\rho''$ du gpe \deleted{\unsure{$H$}} ``abstrait'' $H$ des autom orientés du \deleted{cube} tétraèdre, ou du cube \ldots\ \deleted{\unsure{laissant|laissons} l'icos-}

Notons qu'un autom de $(G,\rho,\rho',\rho'')$ est un autom de $(G,H)$ qui commute à $\rho,\rho',\rho''$ i.e. à $H$, donc \deleted{est} trivial (car le centr. de $H$ est réduit au gpe unité). Donc on a catégorie \uline{rigide} : ces objets sont définis à isom. \uline{unique} près, pour la classe d'isom. donnée.

% archive sheet number: 44
